% !TEX root = ../main.tex
\thispagestyle{acknowledgements}
\section*{Abstract}
\addcontentsline{toc}{section}{Abstract}

On-policy reinforcement learning for contact-rich robotic manipulation requires
extensive interaction. This thesis studies one task: planar pushing of a
T-block to a target pose by a simulated UR5e arm in NVIDIA Isaac Lab. Pushing
mechanics couple translation and rotation, so a goal that specifies a position
and an orientation cannot be split into independent sub-problems. \\

Asymmetric Self-Play (ASP) is an automatic curriculum that succeeds on
goal-reaching problems elsewhere, yet on the pushing task it transfers none of
its within-distribution competence to held-out scenes. A plain single-agent policy
with a potential-based reward reaches 79.3\,\% held-out scene success
($\pm$ 3.5\,pp, five seeds). The self-play subjects solve 0.0\,\% on the
T-block and 7.3\,\% on a rotationally symmetric disc. The mechanism is a
generalisation failure compounded by the coupled multi-objective gate. \\

A within-distribution evaluation shows the confinement: on goals the
goal-proposing agent itself produces, the self-play policy solves 61.4\,\% of
disc scenes and 11.7\,\% of T-block scenes, against 7.3\,\% and 0.0\,\% of
held-out scenes. The imitation term is suppressed in discrete macro-actions
(clip fraction 0.000); an eightfold larger batch leaves the result unchanged.
ASP succeeds in the settings prior work used. The thesis
explains why it does not generalise under a single Graphical Processing Unit
budget. For this class of task, principled reward design produces a larger
improvement than structured training.
